\documentclass{article}
\usepackage[utf8]{inputenc}
\usepackage{amsmath,amssymb}
\usepackage[most]{tcolorbox}
\usepackage{geometry}
\geometry{margin=2.5cm}

\newcommand{\step}[1]{\par\noindent\hangindent=3.5em\hangafter=1%
  \makebox[3.5em][l]{\textbf{#1.}}}
\newcommand{\steptag}[1]{\hfill{\small\textsf{[#1]}}}

\newtcolorbox{proofbox}[1]{
  colback=gray!5, colframe=gray!60,
  fonttitle=\bfseries, title={#1},
  breakable, enhanced,
  left=4pt, right=4pt, top=4pt, bottom=4pt,
  before skip=10pt, after skip=10pt
}

\begin{document}

\begin{proofbox}{There are universal a\_corr>0 and C\_corr<infinity with the...}
\step{1} There are universal a\_corr>0 and C\_corr<infinity with the following property: if B is a finite-dimensional C*-algebra, H a finite-dimensional Hilbert space, and T:B->B(H) is linear, dagger-preserving, has ||T\_n(XY)-T\_n(X)T\_n(Y)||<=a||X||||Y|| and ||T\_n(I)-I||<=a at every n, where 0<=a<=a\_corr, then one unital dagger-homomorphism mu:B->B(H) satisfies ||mu\_n-T\_n||<=C\_corr*a at every n. \steptag{validated} \\[6pt]
\step{1.1} Uniform cb control, exact unitalization, and the norm-one diagonal. Assume 0<=a<=a\_corr<=1. At each n and ||X||=1, dagger preservation, the C*-identity in B(C\textasciicircum{}n tensor H), and the stated multiplicative defect give ||T\_n(X)||\textasciicircum{}2=||T\_n(X)\textasciicircum{}dagger T\_n(X)||<=||T\_n(X\textasciicircum{}dagger X)||+a<=M\_n+a, where M\_n=||T\_n||. Hence M\_n\textasciicircum{}2<=M\_n+a and M\_n<2, uniformly in n. Put e=I-T(I), choose a state phi on B, and set S\_0(x)=T(x)+phi(x)e. Since ||phi||\_cb=1, phi(x\textasciicircum{}dagger)=conj(phi(x)), e=e\textasciicircum{}dagger and ||e||<=a, the same level-one S\_0 is dagger-preserving and exactly unital, ||S\_0-T||\_cb<=a, and ||S\_0||\_cb<3. If G\_T and G\_0 are the amplified defects, direct expansion gives G\_0(x,y)=G\_T(x,y)+phi(xy)e-phi(y)T(x)e-phi(x)eT(y)-phi(x)phi(y)e\textasciicircum{}2, with the entrywise amplified version at every n; therefore its uniform amplified bilinear norm b\_0 is at most a+a+2a+2a+a\textasciicircum{}2<=7a. Set c\_u=7. Finally normalized Haar measure on the compact unitary group of B gives D=int U\textasciicircum{}dagger tensor U dU. In finite dimension D is a finite convex combination sum\_s p\_s U\_s\textasciicircum{}dagger tensor U\_s, with sum\_s ||p\_s U\_s\textasciicircum{}dagger||||U\_s||=1. Haar invariance first for unitary X, and then linearity because every element of a unital C*-algebra is a linear combination of unitaries, gives XD=DX; also multiplication(D)=I. Thus D is a dimension-independent norm-one diagonal. For any linear S and amplified bilinear defect G, w prime(x)=sum\_s S(p\_s U\_s\textasciicircum{}dagger)G(U\_s,x) obeys ||w prime||\_cb<=||S||\_cb||G||. Its amplification is the entrywise amplification of this same level-one map, so all later bounds are uniform in n and in the dimensions. \steptag{validated} \\[4pt]
\step{1.2} One normalized Newton correction. Let S:B->B(H) be exactly unital and dagger-preserving, with ||S||\_cb<=4, and let G(x,y)=S(xy)-S(x)S(y) have uniform amplified bilinear norm b. Use the diagonal D=sum\_s a\_s tensor b\_s from node 1.1, where a\_s=p\_s U\_s\textasciicircum{}dagger and b\_s=U\_s, and define w prime(x)=sum\_s S(a\_s)G(b\_s,x), w doubleprime(x)=w prime(x\textasciicircum{}dagger)\textasciicircum{}dagger, w=(w prime+w doubleprime)/2, and S plus=S+w. Then ||w||\_cb<=4b, w is dagger-preserving, and G(y,I)=G(I,y)=0 implies w(I)=0, so S plus is exactly unital and dagger-preserving. Write delta w(x,y)=S(x)w(y)-w(xy)+w(x)S(y). Associativity gives the exact cocycle identity S(x)G(y,z)-G(xy,z)+G(x,yz)-G(x,y)S(z)=0. Centrality of D and multiplication(D)=I then give delta w prime=G-Q prime, where R=sum\_s G(a\_s,b\_s) and Q prime(x,y)=R G(x,y)+sum\_s G(x,a\_s)G(b\_s,y). Hence ||Q prime||<=2b\textasciicircum{}2 at every amplification. Daggering the identity gives delta w doubleprime=G-Q doubleprime with Q doubleprime(x,y)=Q prime(y\textasciicircum{}dagger,x\textasciicircum{}dagger)\textasciicircum{}dagger and the same bound. Therefore the new defect is G\_\{S plus\}=G-delta w-w(x)w(y)=(Q prime+Q doubleprime)/2-w(x)w(y), so its uniform amplified bilinear norm is at most 2b\textasciicircum{}2+16b\textasciicircum{}2=18b\textasciicircum{}2. Thus, with K\_N=18, ||S plus-S||\_cb<=4b and b plus<=K\_N b\textasciicircum{}2; all formulas are amplifications of the same corrected level-one map. \steptag{validated} \\[4pt]
\step{1.3} Convergence and exactness. Let c\_u=7 and K\_N=18 be supplied by nodes 1.1 and 1.2, and choose a\_corr=min(1/(2 K\_N c\_u),1/(16 c\_u)); then K\_N c\_u a\_corr<=1/2 and 8 c\_u a\_corr<=1/2<1. Starting from S\_0 of node 1.1, recursively apply node 1.2. If b\_k is the uniform amplified bilinear defect of S\_k, then b\_k<=b\_0 2\textasciicircum{}\{-k\}, every iterate has cb norm below 4, and sum\_k ||S\_\{k+1\}-S\_k||\_cb<=8b\_0. Hence S\_k converges in cb norm to one exactly unital dagger-preserving linear map mu. Since b\_k tends to zero, continuity of multiplication gives mu(xy)=mu(x)mu(y). Moreover ||mu-T||\_cb<=||S\_0-T||\_cb+||mu-S\_0||\_cb<=(1+8c\_u)a=57a. Thus C\_corr=57 works, and because convergence is in cb norm, each mu\_n is the amplification of this same level-one mu and ||mu\_n-T\_n||<=57a at every n. \steptag{validated} \\[6pt]
\step{1.3.1} Induction and limit bridge. Invoke the registered validation dependencies 1.1 and 1.2. Let b\_k be the uniform amplified bilinear norm of the multiplicative defect of S\_k. Node 1.1 gives an exactly unital dagger-preserving S\_0 with ||S\_0||\_cb<3, b\_0<=c\_u a, and ||S\_0-T||\_cb<=a. Inductively suppose S\_k is exactly unital and dagger-preserving, ||S\_k||\_cb<4, and b\_k<=b\_0 2\textasciicircum{}\{-k\}. Node 1.2 produces the same level-one map S\_\{k+1\} at all amplifications, still exactly unital and dagger-preserving, with ||S\_\{k+1\}-S\_k||\_cb<=4b\_k and b\_\{k+1\}<=K\_N b\_k\textasciicircum{}2. Since K\_N b\_0<=K\_N c\_u a<=1/2, b\_\{k+1\}<=(K\_N b\_0)b\_k<=b\_k/2<=b\_0 2\textasciicircum{}\{-(k+1)\}. The degenerate-safe geometric estimate in child 1.3.1.1 gives ||S\_k-S\_0||\_cb<=8b\_0<=8c\_u a<1, including b\_0=0, so ||S\_k||\_cb<4 and the induction closes. It also gives sum\_k ||S\_\{k+1\}-S\_k||\_cb<=8b\_0. Therefore S\_k is Cauchy in cb norm and converges to a single level-one linear map mu; exact unitality and dagger preservation pass to the limit. For fixed x,y, the boundedness ||S\_k||\_cb<4, convergence S\_k->mu, and b\_k->0 show S\_k(xy)-S\_k(x)S\_k(y)->mu(xy)-mu(x)mu(y) while the left side has norm at most b\_k||x||||y|$|-\rangle$0. Hence mu is multiplicative. Finally ||mu-T||\_cb<=||S\_0-T||\_cb+sum\_k||S\_\{k+1\}-S\_k||\_cb<=a+8c\_u a=57a. Cb convergence is sup\_n||(S\_k-mu)\_n|$|-\rangle$0, so every mu\_n is the amplification of this same mu and ||mu\_n-T\_n||<=57a. \steptag{validated} \\[6pt]
\step{1.3.1.1} Degenerate-safe geometric estimate. Suppose b\_k<=b\_0 2\textasciicircum{}\{-k\} and ||S\_\{k+1\}-S\_k||\_cb<=4b\_k, as obtained in the parent induction. For every k, ||S\_k-S\_0||\_cb<=sum\_\{j<k\}4b\_j<=4b\_0 sum\_\{j=0\}\textasciicircum{}\{k-1\}2\textasciicircum{}\{-j\}<=8b\_0<=8c\_u a<=8c\_u a\_corr<=1/2<1. This includes b\_0=0: then 0<=b\_k<=0 gives b\_k=0 for all k, and ||S\_\{k+1\}-S\_k||\_cb<=0 forces every correction increment to vanish. In every case, ||S\_k||\_cb<=||S\_0||\_cb+||S\_k-S\_0||\_cb<3+1=4, so the Newton hypothesis closes at the next step without assuming b\_0>0. Likewise sum\_\{k>=0\}||S\_\{k+1\}-S\_k||\_cb<=4 sum\_\{k>=0\}b\_k<=4b\_0 sum\_\{k>=0\}2\textasciicircum{}\{-k\}=8b\_0. These are precisely the finite-displacement and Cauchy estimates used by the parent. \steptag{validated} \\[4pt]
\end{proofbox}

\end{document}
